\chapter{Solución analítica de la frontera eficiente}
\label{ap:frontera-analitica}

Cuando se eliminan las restricciones de desigualdad ($w_i \ge 0$ y
$w_i \le w_{\max}$), el problema de Markowitz admite una solución analítica
cerrada que constituye uno de los resultados clásicos de la teoría de
carteras \cite{Merton1972}. Este apéndice recoge su derivación completa, a
la que se remite desde el Capítulo~\ref{ch:formulacion}. En todo el desarrollo
se supone que $\boldsymbol{\Sigma}$ es \textbf{definida positiva}
($\boldsymbol{\Sigma} \succ 0$), hipótesis bajo la cual $\boldsymbol{\Sigma}^{-1}$
existe; en la práctica se cumple cuando el número de observaciones supera al de
activos y ninguna combinación de activos es redundante (cuestión que se discute
en el Capítulo~\ref{ch:limitaciones}).

\section{Caso sin restricciones de desigualdad}

El problema se reduce a:
\begin{equation}
    \begin{aligned}
        \min_{\mathbf{w}} \quad & \frac{1}{2} \mathbf{w}^T
            \boldsymbol{\Sigma} \mathbf{w} \\
        \text{s.a.} \quad & \mathbf{1}^T \mathbf{w} = 1, \\
        & \mathbf{w}^T \boldsymbol{\mu} = \mu_p.
    \end{aligned}
    \label{eq:qp-sin-desigualdad}
\end{equation}

Nótese que la restricción de rentabilidad se ha escrito como igualdad
($=$) en lugar de desigualdad ($\ge$): en el óptimo, cualquier objetivo
$\mu_p$ alcanzable se saturará porque añadir rentabilidad sin aumentar
riesgo es siempre beneficioso.

\section{Derivación mediante multiplicadores de Lagrange}

El Lagrangiano del problema (\ref{eq:qp-sin-desigualdad}) es:
\begin{equation}
    \mathcal{L}(\mathbf{w}, \lambda, \gamma)
    = \frac{1}{2} \mathbf{w}^T \boldsymbol{\Sigma} \mathbf{w}
    + \lambda (1 - \mathbf{1}^T \mathbf{w})
    + \gamma (\mu_p - \mathbf{w}^T \boldsymbol{\mu}).
    \label{eq:lagrangiano-sin-desigualdad}
\end{equation}

La condición de estacionariedad $\nabla_{\mathbf{w}} \mathcal{L} =
\mathbf{0}$ proporciona:
\begin{equation}
    \boldsymbol{\Sigma} \mathbf{w} - \lambda \mathbf{1}
    - \gamma \boldsymbol{\mu} = \mathbf{0}
    \quad \Longrightarrow \quad
    \mathbf{w} = \boldsymbol{\Sigma}^{-1}
    (\lambda \mathbf{1} + \gamma \boldsymbol{\mu}).
    \label{eq:estacionariedad-analitica}
\end{equation}

Sustituyendo en las restricciones $\mathbf{1}^T \mathbf{w} = 1$ y
$\boldsymbol{\mu}^T \mathbf{w} = \mu_p$, se obtiene un sistema lineal
$2 \times 2$ para $(\lambda, \gamma)$:
\begin{equation}
    \begin{pmatrix}
        A & B \\
        B & C
    \end{pmatrix}
    \begin{pmatrix}
        \lambda \\
        \gamma
    \end{pmatrix}
    = \begin{pmatrix}
        1 \\
        \mu_p
    \end{pmatrix},
    \label{eq:sistema-lambda-gamma}
\end{equation}
donde se han definido las constantes escalares:
\begin{equation}
    A = \mathbf{1}^T \boldsymbol{\Sigma}^{-1} \mathbf{1}, \qquad
    B = \mathbf{1}^T \boldsymbol{\Sigma}^{-1} \boldsymbol{\mu}, \qquad
    C = \boldsymbol{\mu}^T \boldsymbol{\Sigma}^{-1} \boldsymbol{\mu}.
    \label{eq:constantes-abc}
\end{equation}

Estas tres constantes dependen exclusivamente de $\boldsymbol{\Sigma}$
y $\boldsymbol{\mu}$; no dependen del objetivo $\mu_p$. Como
$\boldsymbol{\Sigma} \succ 0$, también $\boldsymbol{\Sigma}^{-1} \succ 0$,
de modo que $\langle \mathbf{x}, \mathbf{y} \rangle :=
\mathbf{x}^T \boldsymbol{\Sigma}^{-1} \mathbf{y}$ define un producto interior.
En esta métrica $A = \langle \mathbf{1}, \mathbf{1} \rangle$,
$C = \langle \boldsymbol{\mu}, \boldsymbol{\mu} \rangle$ y
$B = \langle \mathbf{1}, \boldsymbol{\mu} \rangle$, por lo que la desigualdad
de Cauchy-Schwarz da $B^2 \le AC$, con igualdad si y solo si $\boldsymbol{\mu}$
y $\mathbf{1}$ son linealmente dependientes (es decir, todos los activos tienen
la misma rentabilidad esperada). Bajo la hipótesis de que $\boldsymbol{\mu}$ y
$\mathbf{1}$ son linealmente independientes, se tiene $\Delta = AC - B^2 > 0$,
lo que garantiza que la matriz del sistema~(\ref{eq:sistema-lambda-gamma}) es
invertible.

Resolviendo el sistema:
\begin{equation}
    \lambda = \frac{C - B\mu_p}{\Delta}, \qquad
    \gamma = \frac{A\mu_p - B}{\Delta}.
    \label{eq:lambda-gamma-solucion}
\end{equation}

\section{Ecuación de la frontera eficiente}

Sustituyendo $\lambda$ y $\gamma$ en la expresión de la varianza
$\sigma_p^2 = \mathbf{w}^T \boldsymbol{\Sigma} \mathbf{w}$, se obtiene
la ecuación de la frontera eficiente en el espacio $(\sigma_p, \mu_p)$:
\begin{equation}
    \boxed{
    \sigma_p^2 = \frac{A \mu_p^2 - 2B \mu_p + C}{\Delta}.
    }
    \label{eq:hiperbola}
\end{equation}

Esta ecuación describe una \textbf{hipérbola} en el plano
$(\sigma_p, \mu_p)$ ---o una parábola en el plano $(\sigma_p^2, \mu_p)$---,
cuya rama derecha ($\sigma_p \ge 0$) es la \textbf{frontera de mínima
varianza}: el lugar geométrico de las carteras de menor riesgo para cada
nivel de rentabilidad. Su asíntota superior tiene pendiente
$\sqrt{\Delta / A}$. La \textbf{frontera eficiente} propiamente dicha es solo
la mitad superior de esa rama, $\mu_p \ge \mu_{\text{GMV}} = B/A$ (el vértice
que se deduce en la ecuación~\ref{eq:mu-gmv}): por debajo del vértice, toda
cartera está dominada por otra de igual riesgo y mayor rentabilidad, y por
tanto no es eficiente.

\section{Cartera de mínima varianza global (GMV)}

El vértice de la hipérbola corresponde a la cartera de mínima varianza
global. Derivando (\ref{eq:hiperbola}) respecto a $\mu_p$ e igualando a
cero:
\begin{equation}
    \frac{d\sigma_p^2}{d\mu_p}
    = \frac{2A\mu_p - 2B}{\Delta} = 0
    \quad \Longrightarrow \quad
    \boxed{\mu_{\text{GMV}} = \frac{B}{A}}.
    \label{eq:mu-gmv}
\end{equation}

Sustituyendo en (\ref{eq:lambda-gamma-solucion}) se obtiene $\lambda =
1/A$, $\gamma = 0$, y el vector de pesos de la GMV es:
\begin{equation}
    \boxed{\mathbf{w}_{\text{GMV}}
    = \frac{\boldsymbol{\Sigma}^{-1} \mathbf{1}}
           {\mathbf{1}^T \boldsymbol{\Sigma}^{-1} \mathbf{1}}.
    }
    \label{eq:w-gmv-analitica}
\end{equation}

La varianza de la GMV es $\sigma^2_{\text{GMV}} = 1/A$.

\section{Cartera de máximo ratio de Sharpe}

La cartera tangente (máximo ratio de Sharpe) se obtiene maximizando
$S_p = (\mu_p - r_f) / \sigma_p$ sobre la frontera. Con $r_f = 0$, la
condición de optimalidad es:
\begin{equation}
    \frac{\mu_p}{\sigma_p} = \frac{d\mu_p}{d\sigma_p}
    \quad \Longrightarrow \quad
    \boxed{\mu_T = \frac{C}{B}, \quad
           \sigma_T^2 = \frac{C}{B^2}}.
    \label{eq:max-sharpe-analitica}
\end{equation}

Y los pesos correspondientes:
\begin{equation}
    \boxed{\mathbf{w}_{T}
    = \frac{\boldsymbol{\Sigma}^{-1} \boldsymbol{\mu}}
           {\mathbf{1}^T \boldsymbol{\Sigma}^{-1} \boldsymbol{\mu}}.
    }
    \label{eq:w-tangente-analitica}
\end{equation}

\section{Validación con la implementación numérica}

La solución analítica (\ref{eq:w-gmv-analitica}) proporciona una
referencia exacta contra la que verificar la implementación numérica
con \texttt{cvxpy}. En el test
\texttt{test\_gmv\_coincide\_\allowbreak con\_\allowbreak analitica\_\allowbreak sin\_\allowbreak restricciones}
del código, se genera una matriz $\boldsymbol{\Sigma}$ sintética
(de $3 \times 3$), se resuelve el QP sin restricciones de desigualdad
y se comprueba que la solución numérica coincide con la expresión
(\ref{eq:w-gmv-analitica}) con un error relativo inferior a $10^{-5}$.
Este acuerdo valida tanto la formulación del QP como su resolución
numérica.
